%%%%%%%%%%%%%%%%%%%%%%%%%%%%%%%%%%%%%%%%%%%%%%%%%%%%%%%%%%%%%%%%%%%%%%%%%%%%%%%%%%
\begin{frame}[fragile]\frametitle{}

\begin{center}
{\Large OpenCode}
\end{center}
\end{frame}

%%%%%%%%%%%%%%%%%%%%%%%%%%%%%%%%%%%%%%%%%%%%%%%%%%%%%%%%%%%%%%%%%%%%%%%%%%%%%%%%%%
\begin{frame}[fragile]\frametitle{OpenCode:  Rapid Growth}
\begin{columns}
    \begin{column}[T]{0.6\linewidth}
  \begin{itemize}
    \item \textbf{90K+ GitHub stars} accumulated in less than a year:  one of the fastest-growing developer tools in history.
    \item \textbf{650,000+ monthly users} as of early 2026.
    \item \textbf{Official partnerships:} OpenAI and GitHub have formally partnered with OpenCode.
    \item Born as a Claude Code alternative; evolved into a model-agnostic agent arguably more capable than what it was imitating.
    \item Even the creator of Supabase publicly praised its momentum.
    \item Developers who were paying \$20/month for Claude Code realized they could use any model from any provider, from their terminal, without restrictions.
  \end{itemize}

    \end{column}
    \begin{column}[T]{0.4\linewidth}
\begin{center}
\includegraphics[width=\linewidth,keepaspectratio]{opencode4}
\end{center}

{\tiny (Ref: JP Caparas, ``The definitive guide to OpenCode: from first install to production workflows'', Dev Genius / reading.sh, Feb 2026)}
    \end{column}
  \end{columns}
\end{frame}

%%%%%%%%%%%%%%%%%%%%%%%%%%%%%%%%%%%%%%%%%%%%%%%%%%%%%%%%%%%%%%%%%%%%%%%%%%%%%%%%%%
\begin{frame}[fragile]\frametitle{OpenCode Ecosystem}
\begin{center}
\includegraphics[width=\linewidth,keepaspectratio]{opencode3}
\end{center}

{\tiny (Ref: OpenCode CRASH Course - Data Science Basics)}
\end{frame}


%%%%%%%%%%%%%%%%%%%%%%%%%%%%%%%%%%%%%%%%%%%%%%%%%%%%%%%%%%%%%%%%%%%%%%%%%%%%%%%%%%
\begin{frame}[fragile]\frametitle{What is OpenCode?}
  \begin{itemize}
    \item Open-source coding agent built for the terminal (TUI).
	\item \textbf{IDE Integration possible}:	Works with VS Code, Cursor, and any terminal-supporting IDE
	\item \textbf{GitHub Integration}: Automate issues and PRs directly from comments
    \item \textbf{Model-Agnostic:} Supports 75+ AI models (OpenAI, Claude, Groq, Ollama, and more).
      \begin{itemize}
        \item \textit{Local:} Ollama (fully offline, zero data egress)
        \item \textit{Cloud:} OpenAI, Claude, Groq, Gemini, etc.
      \end{itemize}	
    \item \textbf{Privacy-Focused:} Runs locally or uses your own API keys, your code stays on your machine.
    \item \textbf{Built With:} Rust/Go for performance; Bubble Tea for the TUI layer.
  \end{itemize}
\end{frame}

%%%%%%%%%%%%%%%%%%%%%%%%%%%%%%%%%%%%%%%%%%%%%%%%%%%%%%%%%%%%%%%%%%%%%%%%%%%%%%%%%%
\begin{frame}[fragile]\frametitle{Capabilities}
Think of it as having a senior developer sitting next to you
  \begin{itemize}
    \item Writing and debugging code
    \item Understanding existing codebases
    \item Refactoring and optimization
    \item Answering technical questions
    \item Automating repetitive tasks
  \end{itemize}
\end{frame}


%%%%%%%%%%%%%%%%%%%%%%%%%%%%%%%%%%%%%%%%%%%%%%%%%%%%%%%%%%%%%%%%%%%%%%%%%%%%%%%%%%
\begin{frame}[fragile]\frametitle{Key Features}
  \begin{itemize}
    \item \textbf{Interactive TUI:} Smooth terminal UI built with Bubble Tea; keyboard-driven workflow.
    \item \textbf{Contextual Understanding:} Reads your entire repo, understands dependencies, structure, languages.
    \item \textbf{Agentic Workflows:} Plans changes, searches files, executes shell commands, modifies code autonomously.
    \item \textbf{Plan / Build Modes/Agents:} Review proposed changes before applying,  deterministic control.
    \item \textbf{Slash Commands:} \texttt{/init}, \texttt{/help}, \texttt{/models}, \texttt{/share}, and more.
    \item \textbf{Context Mentions:} \texttt{@file:}, \texttt{@workspace} for scoped AI queries.
  \end{itemize}
\end{frame}

%%%%%%%%%%%%%%%%%%%%%%%%%%%%%%%%%%%%%%%%%%%%%%%%%%%%%%%%%%%%%%%%%%%%%%%%%%%%%%%%%%
\begin{frame}[fragile]\frametitle{Core Integration}
  \begin{itemize}
    \item \textbf{LSP Integration:} Language Server Protocol for real-time code intelligence (completions, diagnostics, go-to-def).
    \item \textbf{MCP (Model Context Protocol):} Plug in external tools and data sources to extend agent capabilities.
    \item \textbf{ACP Support:} Agent Communication Protocol for inter-agent coordination.
    \item \textbf{Agent Skills \& Custom Tools:} Define reusable agent behaviors and wire in custom shell/API tools.
    \item \textbf{Rules \& Permissions:} Fine-grained control over what the agent can read, write, or execute.
    \item \textbf{Multi-Provider Routing:} Failover and cost-based routing across model providers.
  \end{itemize}
\end{frame}



%%%%%%%%%%%%%%%%%%%%%%%%%%%%%%%%%%%%%%%%%%%%%%%%%%%%%%%%%%%%%%%%%%%%%%%%%%%%%%%%%%
\begin{frame}[fragile]\frametitle{OpenCode vs Other Tools}
\begin{center}
\includegraphics[width=\linewidth,keepaspectratio]{opencode1}
\end{center}

{\tiny (Ref: OpenCode - nxcode.io}
\end{frame}



%%%%%%%%%%%%%%%%%%%%%%%%%%%%%%%%%%%%%%%%%%%%%%%%%%%%%%%%%%%%%%%%%%%%%%%%%%%%%%%%%%
\begin{frame}[fragile]\frametitle{Installation:  All Platforms}
  \begin{itemize}
	\item \textbf{Desktop App:} Download directly from opencode.ai/download for a native desktop experience.
    \item \textbf{Quick Install (Linux/macOS):} \lstinline!$ curl -fsSL https://opencode.ai/install | sh!
	\item \textbf{npm/bun:} \lstinline|npm i -g opencode-ai@latest| or \lstinline|bun add -g opencode-ai|
    \item \textbf{macOS (Homebrew):} \lstinline|brew install sst/tap/opencode|
    \item \textbf{Windows (Chocolatey):} \lstinline|choco install opencode|
    \item \textbf{Windows (Scoop):} \lstinline|scoop install opencode|
    \item \textbf{Arch Linux (paru):} \lstinline|paru -S opencode-bin|
    \item Verify installation: \lstinline|opencode --version|
	\item If installed using npm on Windows 11 then 
		\begin{itemize}
			\item Installation directory: \lstinline|%USERPROFILE%\AppData\Roaming\npm\node_modules\opencode-ai\|
			
			\item Configuration and data (session data, logs, and authentication tokens): 
			\lstinline|%USERPROFILE%\.local\share\opencode\auth.json|
			
			\item Global plugins: 
			\lstinline|%USERPROFILE%\.config\opencode\plugins\|
			
			\item Global configuration file for user-wide settings (providers, models, etc.): 
			\lstinline|%USERPROFILE%\.config\opencode\opencode.json|
		\end{itemize}
  \end{itemize}
\end{frame}

%%%%%%%%%%%%%%%%%%%%%%%%%%%%%%%%%%%%%%%%%%%%%%%%%%%%%%%%%%%%%%%%%%%%%%%%%%%%%%%%%%
\begin{frame}[fragile]\frametitle{Getting Started:  Connecting a Provider}

\begin{columns}
    \begin{column}[T]{0.6\linewidth}
  \begin{itemize}
    \item \textbf{Launch:} Run \texttt{opencode} in any project directory. The TUI loads; your first step is connecting a provider.
    \item \textbf{Connect:} Type \texttt{/connect} → interactive menu → select provider → paste API key.
    \item \textbf{Bring Your Own Keys:} Set \texttt{OPENAI\_API\_KEY}, \texttt{ANTHROPIC\_API\_KEY}, or point to Ollama.
    \item \textbf{Configure Provider:} Global config \lstinline|~/.config/opencode/opencode.json| or project-specific \lstinline|./opencode.json|
    \item \textbf{Config Format:} JSON-based \texttt{opencode.json}; supports project-level and global overrides.
    \item \textbf{First Step After Connect:} Run \texttt{/init} to ground the agent in your project:  scans the codebase and creates \texttt{AGENTS.md}.
  \end{itemize}

    \end{column}
    \begin{column}[T]{0.4\linewidth}
\begin{center}
\includegraphics[width=\linewidth,keepaspectratio]{opencode5}
\end{center}

{\tiny (Ref: JP Caparas, ``The definitive guide to OpenCode: from first install to production workflows'', Dev Genius / reading.sh, Feb 2026)}
    \end{column}
  \end{columns}
  
\end{frame}

%%%%%%%%%%%%%%%%%%%%%%%%%%%%%%%%%%%%%%%%%%%%%%%%%%%%%%%%%%%%%%%%%%%%%%%%%%%%%%%%%%
\begin{frame}[fragile]\frametitle{Providers in OpenCode}

\textbf{OpenCode supports 75+ LLM providers.} Here are the most useful options.

\vspace{0.3cm}
\textbf{OpenCode Zen (Curated Gateway)}

\begin{itemize}
    \item Pre-benchmarked models with optimized configurations
    \item No setup headaches — models just work
    \item Access to free models (GLM 4.7 Free, Kimi K2.5 Free, Big Pickle)
    \item Competitive pricing (often cheaper than direct APIs)
    \item Team management features built-in
\end{itemize}

\vspace{0.3cm}
\textbf{Pricing (per 1M tokens):}
\begin{itemize}
    \item Claude Sonnet 4.5: \$3.00 input / \$15.00 output
    \item Qwen3 Coder 480B: \$0.45 input / \$1.50 output
    \item Free models: \$0 (may train on user data)
\end{itemize}

\vspace{0.3cm}
\textbf{Setup:}
\begin{itemize}
    \item Sign in at \texttt{opencode.ai/auth}
    \item Add billing details
    \item Run \texttt{/connect}, select OpenCode, paste your key
\end{itemize}

\end{frame}

%%%%%%%%%%%%%%%%%%%%%%%%%%%%%%%%%%%%%%%%%%%%%%%%%%%%%%%%%%%%%%%%%%%%%%%%%%%%%%%%%%
\begin{frame}[fragile]\frametitle{Custom OpenAI-Compatible Providers}

\textbf{Connection Steps:}
\begin{itemize}
    \item Run \texttt{/connect}
    \item Select \texttt{Other}
    \item Enter provider ID
    \item Paste your API key
\end{itemize}

\textbf{Use any provider with an OpenAI-compatible API.}

\begin{lstlisting}
{
  "$schema": "https://opencode.ai/config.json",
  "provider": {
    "my-provider": {
      "npm": "@ai-sdk/openai-compatible",
      "name": "My Custom Provider",
      "options": {
        "baseURL": "https://api.my-provider.com/v1"
      },
      "models": {
        "my-model": {
          "name": "My Model" }}}}}
\end{lstlisting}



\end{frame}

%%%%%%%%%%%%%%%%%%%%%%%%%%%%%%%%%%%%%%%%%%%%%%%%%%%%%%%%%%%%%%%%%%%%%%%%%%%%%%%%%%
\begin{frame}[fragile]\frametitle{Plan vs Build Mode}

Switch between them using the Tab key.

Plan Mode:
  \begin{itemize}
    \item Read-only - Cannot make changes to your code
    \item Analyzes and explores your codebase
    \item Suggests implementation strategies
    \item Perfect for understanding before modifying
  \end{itemize}
  
Build Mode (Default)
  \begin{itemize}
    \item Full access - Can read, write, and modify files
    \item Executes code changes
    \item Runs commands
    \item Best for actual development work
  \end{itemize}  
  
Pro Tip: Always start complex features in Plan mode to outline your approach, then switch to Build mode for implementation.  
\end{frame}

%%%%%%%%%%%%%%%%%%%%%%%%%%%%%%%%%%%%%%%%%%%%%%%%%%%%%%%%%%%%%%%%%%%%%%%%%%%%%%%%%%
\begin{frame}[fragile]\frametitle{Context Mentions \& Image Support}
  \begin{itemize}
    \item \textbf{\texttt{@file:path}} --- focus the agent on a specific file:
      \lstinline|@file:main.go fix the bug where user input isn't sanitized|
    \item \textbf{\texttt{@workspace}} --- ask about the broader codebase:
      \lstinline|Show me how API routes are structured|
    \item \textbf{Onboarding power:} Instead of manually \texttt{grep}-ing and reading files, just ask. Especially useful when exploring an unfamiliar codebase.
    \item \textbf{Image drag-and-drop:} Drag and drop screenshots or mockups directly into the terminal. OpenCode scans them and includes them in context --- useful for implementing UI from a design mockup.
    \item \textbf{Tab to Plan mode:} For complex changes, use \texttt{Tab} to switch to Plan mode, describe the feature, and review the proposed approach before any code is written.
  \end{itemize}

\end{frame}

%%%%%%%%%%%%%%%%%%%%%%%%%%%%%%%%%%%%%%%%%%%%%%%%%%%%%%%%%%%%%%%%%%%%%%%%%%%%%%%%%%
\begin{frame}[fragile]\frametitle{Session Management}
  \begin{itemize}
    \item \textbf{\texttt{/sessions}} (aliases: \texttt{/resume}, \texttt{/continue}) --- opens a session picker showing recent conversations; select one to resume with full context.
    \item \textbf{CLI continue flags:}
      \begin{itemize}
        \item \lstinline|opencode --continue| or \lstinline|opencode -c| --- resumes the most recent session.
        \item \lstinline|opencode session list| --- lists all sessions with their IDs.
        \item \lstinline|opencode session continue <id>| --- resumes a specific session by ID.
      \end{itemize}
    \item \textbf{Session titles:} Auto-generated from your first prompt. Name them upfront for easy retrieval:
      \lstinline|opencode run --title "Payment API refactor" "Refactor the payment processing module"|
    \item \textbf{\texttt{/share}} --- generates a shareable URL for async code review, debugging docs, or team collaboration.
    \item \textbf{\texttt{/compact}} (alias: \texttt{/summarize}) --- compresses the current session context to save tokens while preserving key decisions.
  \end{itemize}
{\tiny (Ref: JP Caparas, ``The definitive guide to OpenCode'', Dev Genius / reading.sh, Feb 2026)}
\end{frame}

%%%%%%%%%%%%%%%%%%%%%%%%%%%%%%%%%%%%%%%%%%%%%%%%%%%%%%%%%%%%%%%%%%%%%%%%%%%%%%%%%%
\begin{frame}[fragile]\frametitle{Undo, Redo \& Non-Interactive Mode}
  \begin{itemize}
    \item \textbf{\texttt{/undo}} --- reverts file changes from the last agent response; restores original prompt so you can rephrase and retry.
    \item \textbf{\texttt{/redo}} --- re-applies changes that were undone.
    \item \textbf{Non-interactive mode} (\texttt{opencode run}) --- pipe prompts for scripting and CI:

    \item \textbf{GitHub Action:} Use \texttt{opencode run} in CI pipelines --- e.g., auto-respond to PR comments, generate changelogs, run AI-assisted linting.
  \end{itemize}
  
\begin{lstlisting}
# One-shot task
opencode run "Add JSDoc comments to all exported functions in src/"

# With title for easy retrieval
opencode run --title "Add JSDoc" "Add JSDoc to src/"

# Continue last session non-interactively
opencode -c run "Now add unit tests for those functions"
\end{lstlisting}
  
\end{frame}

%%%%%%%%%%%%%%%%%%%%%%%%%%%%%%%%%%%%%%%%%%%%%%%%%%%%%%%%%%%%%%%%%%%%%%%%%%%%%%%%%%
\begin{frame}[fragile]\frametitle{}

\begin{center}
{\Large Configuration}
\end{center}
\end{frame}


%%%%%%%%%%%%%%%%%%%%%%%%%%%%%%%%%%%%%%%%%%%%%%%%%%%%%%%%%%%%%%%%%%%%%%%%%%%%%%%%%%
\begin{frame}[fragile]\frametitle{Overview}
  \begin{itemize}
    \item OpenCode is fully configurable at \textbf{two scopes}: Global (\texttt{\textasciitilde/.config/opencode/}) and Project (\texttt{.opencode/} in repo root).
    \item \textbf{Config priority} (highest → lowest): CLI flags → Custom env (\texttt{OPENCODE\_CONFIG\_DIR}) → Project \texttt{opencode.json} → Global \texttt{opencode.json} → Remote org defaults (\texttt{.well-known/opencode}).
    \item \textbf{Subdirectories} (plural names preferred): \texttt{agents/}, \texttt{skills/}, \texttt{tools/}, \texttt{modes/}, \texttt{commands/}, \texttt{themes/}, \texttt{plugins/}.
    \item \textbf{Schema-validated:} Config is validated against \texttt{opencode.ai/config.json}:  editors autocomplete fields automatically.
    \item \textbf{Claude Code compatibility:} Falls back to \texttt{CLAUDE.md} / \texttt{\textasciitilde/.claude/} conventions if OpenCode files are absent.
    \item Each section below covers: what ships built-in → what you can customize → how to create your own.
  \end{itemize}
  
\textbf{The core problem:} Context is a finite budget. Dumping everything into one file degrades quality:  the agent buries critical constraints under noise. OpenCode solves this with \textit{tiered, purposeful} context loading.  
\end{frame}

%%%%%%%%%%%%%%%%%%%%%%%%%%%%%%%%%%%%%%%%%%%%%%%%%%%%%%%%%%%
%% OpenCode:  Why So Many Customization Layers?
%%%%%%%%%%%%%%%%%%%%%%%%%%%%%%%%%%%%%%%%%%%%%%%%%%%%%%%%%%%

%%%%%%%%%%%%%%%%%%%%%%%%%%%%%%%%%%%%%%%%%%%%%%%%%%%%%%%%%%%%%%%%%%%%%%%%%%%%%%%%%%
\begin{frame}[fragile]\frametitle{Customization Layers:  The Mental Model}
  \begin{itemize}

    \item \textbf{AGENTS.md:  ``The Onboarding Doc'':} \textit{Always-on}, repo-committed, team-shared. Tells \textit{any} agent how this project works: stack, build commands, test steps, conventions. Think of it as the README for AI:  written once, applied everywhere.
    \item \textbf{Rules:  ``Standing Orders'':} \textit{Always-on}, injected into every session. Global rules (\texttt{\textasciitilde/.config/opencode/AGENTS.md}) encode personal preferences never shared with the team. Multi-file rules via \texttt{opencode.json} allow conditional loading per glob/path. Same intent as AGENTS.md:  just scoped differently (personal vs. project vs. package).
	\item OpenCode has three extension mechanisms: commands, skills, and agents. Commands are what to do. Skills are how to do it. Agents are who does it. 
	\item \item \textbf{Commands:  ``The shortcuts for repetitive prompts'':} stored in \lstinline|Create .opencode/commands/test.md| and triggered with $/$
    \item \textbf{Skills:  ``The On-Call Expert'':} \textit{Lazy-loaded}: Unlike commands (which execute), skills inform. The agent reads only the name + description at startup, then fetches full content only when the task is relevant. Prevents token waste on instructions that don't apply to the current task. Think: nobody memorizes the entire wiki before starting a ticket.
    \item \textbf{Agents:  ``The Specialist Persona'':} Own model, temperature, tool set, permissions, and system prompt. Invoked by \texttt{@mention} or as a subagent. Not just instructions:  a \textit{different actor} with a different role (e.g., a read-only \texttt{Review} agent vs. a write-enabled \texttt{Build} agent).
  \end{itemize}
\end{frame}

%%%%%%%%%%%%%%%%%%%%%%%%%%%%%%%%%%%%%%%%%%%%%%%%%%%%%%%%%%%%%%%%%%%%%%%%%%%%%%%%%%
\begin{frame}[fragile]\frametitle{Customization Layers:  When to Use What}

\textbf{The layering principle:} Rules = what to always know. Skills = what to know when needed. Agents = who is doing the work and with what capabilities.

  \begin{itemize}
    \item \textbf{Use AGENTS.md when:} The instruction is \textit{project-wide}, stable, and useful to every agent, every session:  build steps, test commands, folder conventions, PR rules. Commit it to Git; share with the team. Auto-generate with \texttt{/init}.
    \item \textbf{Use Rules when:} The instruction is always relevant but \textit{not} project-specific:  personal tone, preferred libraries, output format, privacy rules. Global rules stay on your machine, never in the repo.
    \item \textbf{Use Skills when:} The instruction is \textit{task-specific} and only occasionally needed:  e.g., a \texttt{git-release} skill for changelogs, a \texttt{db-migration} skill for schema changes. Keeps context lean; agent loads it only when it matches the task.
    \item \textbf{Use Custom Agents when:} You need a \textit{different set of capabilities} for a role:  a \texttt{Review} agent that can't write files, a \texttt{Debug} agent with only bash + read, a \texttt{Security} agent using a different model.
    \item \textbf{Use Commands (\texttt{/}) when:} You have a \textit{repetitive prompt}:  a parameterized template you type often. Commands invoke agents, pass arguments, and save keystrokes.
  \end{itemize}
\end{frame}

%%%%%%%%%%%%%%%%%%%%%%%%%%%%%%%%%%%%%%%%%%%%%%%%%%%%%%%%%%%%%%%%%%%%%%%%%%%%%%%%%%
\begin{frame}[fragile]\frametitle{Customization Layers:  Summary Table}

      \begin{tabular}{|l|l|l|l|}
        \hline
        \textbf{Layer} & \textbf{Loaded} & \textbf{Scope}  \\
        \hline
        AGENTS.md & Always & Project / Package  \\
        Rules & Always & Personal / Team  \\
        Skills & On-demand & Any  \\
        Agents & When invoked & Session role \\
        Commands & When typed & Shortcut  \\
        \hline
      \end{tabular}

\end{frame}


%%%%%%%%%%%%%%%%%%%%%%%%%%%%%%%%%%%%%%%%%%%%%%%%%%%%%%%%%%%%%%%%%%%%%%%%%%%%%%%%%%
\begin{frame}[fragile]\frametitle{AGENTS.md}
AGENTS.md helps OpenCode understand: project structure, Coding conventions, Technology stack etc

\begin{lstlisting}
## Project: My App

## Tech Stack
- Next.js 14 with App Router
:

## Conventions
- Use functional components
:

## Structure
- /app - Next.js app router pages
- /components - Reusable UI components
:
\end{lstlisting}
\end{frame}

%%%%%%%%%%%%%%%%%%%%%%%%%%%%%%%%%%%%%%%%%%%%%%%%%%%%%%%%%%%%%%%%%%%%%%%%%%%%%%%%%%
\begin{frame}[fragile]\frametitle{Rules:  Persistent AI Instructions}
  \begin{itemize}
    \item \textbf{What:} Plain-text instructions injected into every LLM context. Think of them as standing orders:  coding style, conventions, constraints.
    \item \textbf{Built-in bootstrap:} \texttt{/init} auto-generates \texttt{AGENTS.md} by scanning your repo:  languages, deps, folder structure, patterns.
    \item \textbf{Scope:}
      \begin{itemize}
        \item \textit{Project rules}: \texttt{./AGENTS.md}:  commit to Git; shared across team.
        \item \textit{Global rules}: \texttt{\textasciitilde/.config/opencode/AGENTS.md}:  personal, not committed.
        \item \textit{Claude Code fallbacks}: \texttt{\textasciitilde/.claude/CLAUDE.md} if no \texttt{AGENTS.md} 
      \end{itemize}
    \item \textbf{Multi-file rules} via \texttt{opencode.json}: glob patterns, remote URLs supported.

    \item \textbf{Lazy loading:} Reference \texttt{@rules/style.md} in \texttt{AGENTS.md}:  agent reads files on-demand, not upfront (saves tokens).
    \item \textbf{Create custom:} Write any Markdown file → list it under \texttt{instructions} in \texttt{opencode.json}. No special syntax:  plain instructions work.
  \end{itemize}
  
\begin{lstlisting}
{ "instructions": ["docs/standards.md",
    "test/guidelines.md", "packages/*/AGENTS.md",
    "https://myorg.com/team-rules.md"] }
\end{lstlisting}  
\end{frame}

%%%%%%%%%%%%%%%%%%%%%%%%%%%%%%%%%%%%%%%%%%%%%%%%%%%%%%%%%%%%%%%%%%%%%%%%%%%%%%%%%%
\begin{frame}[fragile]\frametitle{Agents:  Built-in \& Custom}
  \begin{itemize}
    \item \textbf{Two types:} \textit{Primary agents} (you interact with directly, cycle via \texttt{Tab}) and \textit{Subagents} (specialists invoked by primary agents or \texttt{@mention}).
    \item \textbf{Built-in Primary Agents:}
      \begin{itemize}
        \item \texttt{Build}:  default; all tools enabled; full file + shell access.
        \item \texttt{Plan}:  read-only; all mutating tools set to \texttt{ask}; safe for analysis.
      \end{itemize}
    \item \textbf{Built-in Subagents:}
      \begin{itemize}
        \item \texttt{General}:  full tool access; multi-step parallel tasks.
        \item \texttt{Explore}:  read-only codebase navigator; fast file/pattern search.
      \end{itemize}
    \item \textbf{Common custom agents:} \texttt{Review} (read-only + docs), \texttt{Debug} (bash + read), \texttt{Docs} (write, no shell), \texttt{Security} (audit mode).
    \item \textbf{Create custom:} Run \texttt{opencode agents create} → CLI prompts for scope, description, tools, model → writes \texttt{.opencode/agents/<name>.md}.
    \item \textbf{Agent frontmatter} fields: \texttt{model}, \texttt{temperature}, \texttt{mode} (\texttt{primary}/\texttt{subagent}/\texttt{all}), \texttt{tools}, \texttt{permission}, \texttt{steps} (max iterations), \texttt{color}.
    \item \textbf{Set default agent} in \texttt{opencode.json}: \texttt{"default\_agent": "build"}:  applies across TUI, CLI, IDE, GitHub Action.
  \end{itemize}
\end{frame}

%%%%%%%%%%%%%%%%%%%%%%%%%%%%%%%%%%%%%%%%%%%%%%%%%%%%%%%%%%%%%%%%%%%%%%%%%%%%%%%%%%
\begin{frame}[fragile]\frametitle{Tools:  Built-in \& Custom}
  \begin{itemize}
    \item \textbf{What:} Functions the LLM calls to act on your codebase. All built-in tools are enabled by default.
    \item \textbf{Some Built-in Tools:}
      \begin{itemize}
        \item \texttt{bash}:  execute shell commands (\texttt{git}, \texttt{npm}, \texttt{make}, etc.)
        % \item \texttt{edit / multiedit}:  exact string replacement in files
        \item \texttt{write}:  create/overwrite files
        % \item \texttt{read}:  read file contents (supports line ranges)
        % \item \texttt{grep}:  regex search across codebase
        % \item \texttt{glob}:  find files by pattern (\texttt{**/*.ts})
        \item \texttt{ls}:  list directory contents
        % \item \texttt{patch}:  apply unified diff patches
        % \item \texttt{webfetch}:  fetch URLs during sessions
        \item \texttt{skill}:  load a \texttt{SKILL.md} on-demand
        \item \texttt{todo}:  manage task lists for multi-step work
        % \item \texttt{lsp}:  code intelligence (defs, refs, hover) [experimental]
      \end{itemize}
    \item \textbf{Custom Tools:} TypeScript/JS files in \texttt{.opencode/tools/} (project) or \texttt{\textasciitilde/.config/opencode/tools/} (global). Filename = tool name. Can wrap scripts in any language (Python, Bash, Go, etc.).
    \item \textbf{Override built-ins:} Custom tool with same name replaces the built-in (e.g., \texttt{bash.ts} → sandboxed shell).
  \end{itemize}
  
\begin{lstlisting}
// .opencode/tools/db-query.ts
import { tool } from "@opencode-ai/plugin"
export default tool({ description: "Run SQL query",
  args: { query: tool.schema.string() },
  async execute(args) { /* call your DB */ }
})
\end{lstlisting}  
\end{frame}

%%%%%%%%%%%%%%%%%%%%%%%%%%%%%%%%%%%%%%%%%%%%%%%%%%%%%%%%%%%%%%%%%%%%%%%%%%%%%%%%%%
\begin{frame}[fragile]\frametitle{Agent Skills:  Reusable Behavior Packages}
  \begin{itemize}
    \item \textbf{What:} Markdown instruction bundles loaded on-demand by agents via the \texttt{skill} tool. Unlike rules (always loaded), skills are lazy:  the agent picks them when relevant.
    \item \textbf{No built-in skills ship with OpenCode}:  skills are community/team-created. Claude Code-compatible skills from \texttt{\textasciitilde/.claude/skills/} work as fallbacks.
    \item \textbf{Discovery locations} (all checked at startup):
      \begin{itemize}
        \item \texttt{.opencode/skills/<name>/SKILL.md}:  project-local
        \item \texttt{\textasciitilde/.config/opencode/skills/<name>/SKILL.md}:  global
        \item \texttt{.claude/skills/} and \texttt{.agents/skills/}:  compatibility paths
      \end{itemize}
    \item \textbf{Name rules:} Lowercase alphanumeric + single hyphens; 1–64 chars; must match directory name.
    \item \textbf{Permissions:} Control skill access per agent with \texttt{allow / deny / ask} and wildcard patterns (\texttt{internal-*}).
  \end{itemize}		


\end{frame}


%%%%%%%%%%%%%%%%%%%%%%%%%%%%%%%%%%%%%%%%%%%%%%%%%%%%%%%%%%%%%%%%%%%%%%%%%%%%%%%%%%
\begin{frame}[fragile]\frametitle{Agent Skills:  Create a custom skill}
	
\begin{lstlisting}
# .opencode/skills/git-release/SKILL.md
---
name: git-release
description: Create consistent releases and changelogs
license: MIT
compatibility: opencode
---
## What I do
- Draft release notes from merged PRs
- Propose a version bump
- Generate a `gh release create` command
## When to use me
Use when preparing a tagged release.
\end{lstlisting}

\end{frame}

%%%%%%%%%%%%%%%%%%%%%%%%%%%%%%%%%%%%%%%%%%%%%%%%%%%%%%%%%%%%%%%%%%%%%%%%%%%%%%%%%%
\begin{frame}[fragile]\frametitle{Permissions:  Granular Action Control}
  \begin{itemize}
    \item \textbf{What:} Per-tool access policy:  \texttt{allow} (auto-run), \texttt{ask} (prompt user), or \texttt{deny} (block). Defaults: most tools \texttt{allow}; \texttt{doom\_loop} and \texttt{external\_directory} default to \texttt{ask}; \texttt{.env} files denied by default.
    \item \textbf{Scope:} Set globally in \texttt{opencode.json}; override per-agent in agent frontmatter. Agent rules take precedence.

    % \item \textbf{Runtime \texttt{always} approval:} During a session, approve a pattern once → auto-approved for the rest of the session.
    % \item \textbf{External directory guard:} Use \texttt{external\_directory} permission to allow/deny tool calls touching paths outside the working directory.
    \item \textbf{Skill permissions:} grant/deny specific skills by name or pattern.
	\item \textbf{Wildcard global policy + overrides} (last matching rule wins):

  \end{itemize}
  
\begin{lstlisting}
{ "permission": {
    "*": "ask",          // ask for everything by default
    "bash": {
      "*": "ask",
      "git *": "allow",  // git commands auto-approved
      "npm *": "allow",
      "rm *": "deny"     // destructive: always blocked
    },
    "mymcp_*": "ask"   // all MCP tool calls: prompt first
  }
}
\end{lstlisting}  
\end{frame}

%%%%%%%%%%%%%%%%%%%%%%%%%%%%%%%%%%%%%%%%%%%%%%%%%%%%%%%%%%%%%%%%%%%%%%%%%%%%%%%%%%
\begin{frame}[fragile]\frametitle{Modes:  Contextual Tool Profiles}
  \begin{itemize}
    \item \textbf{What:} Named configurations combining a system prompt + tool subset + temperature. Switch modes to change agent behavior without switching agents.
    \item \textbf{Difference from Agents:} Agents are full AI configurations with their own identity and model; Modes are overlays on top of an agent.
    \item \textbf{Built-in Modes:}
      \begin{itemize}
        \item \texttt{Build} (default):  all tools on, full write access
        \item \texttt{Plan}:  analysis only; mutating tools set to \texttt{ask}
      \end{itemize}
    \item \textbf{Useful custom modes:} \texttt{Review} (read-only + docs), \texttt{Debug} (bash + read, no writes), \texttt{Refactor} (edit only, no shell), \texttt{Docs} (write files, no system commands).
    \item \textbf{Create custom mode}:  Markdown file in \texttt{.opencode/modes/} or \texttt{\textasciitilde/.config/opencode/modes/}:

    \item \textbf{Switch modes} in TUI: \texttt{Tab} key cycles between \texttt{Build} and \texttt{Plan} by default.
  \end{itemize}
  

\end{frame}

%%%%%%%%%%%%%%%%%%%%%%%%%%%%%%%%%%%%%%%%%%%%%%%%%%%%%%%%%%%%%%%%%%%%%%%%%%%%%%%%%%
\begin{frame}[fragile]\frametitle{Modes:  Custom}
\begin{lstlisting}
---
name: debug
description: Investigate issues, no file changes
tools:
  bash: true
  read: true
  grep: true
  write: false
  edit: false
temperature: 0.2
---
You are in debug mode. Diagnose and report only.
Do NOT modify any files. Use bash and read tools
to investigate. Output a structured findings report.
\end{lstlisting}  
\end{frame}
%%%%%%%%%%%%%%%%%%%%%%%%%%%%%%%%%%%%%%%%%%%%%%%%%%%%%%%%%%%%%%%%%%%%%%%%%%%%%%%%%%
\begin{frame}[fragile]\frametitle{Commands:  Slash-Invokable Shortcuts}
  \begin{itemize}
    \item \textbf{What:} Reusable prompt templates invoked with \texttt{/} in the TUI. Eliminate repetitive prompt writing for common workflows.
    \item \textbf{Built-in commands:}
      \begin{itemize}
        \item \texttt{/init}:  scan project, generate \texttt{AGENTS.md}
        \item \texttt{/models}:  browse and switch active model
        % \item \texttt{/connect}:  authenticate a new provider
        % \item \texttt{/share}:  generate a shareable session URL
        \item \texttt{/help}:  list all available commands
      \end{itemize}
    \item \textbf{Create custom commands} in \texttt{opencode.json} (or as \texttt{.opencode/commands/<name>.md} files):

    \item \textbf{\texttt{\$ARGUMENTS}:} Anything typed after \texttt{/command-name} is injected here:  makes commands parameterized.
    \item \textbf{Team sharing:} Commit \texttt{.opencode/commands/} to Git:  team gets it.
  \end{itemize}
  
\begin{lstlisting}
{ "command": {
    "test-report": {
      "description": "Full test suite with coverage",
      "template": "Run test suite with coverage report. Suggest fixes."
    },
  }
}
\end{lstlisting}  
\end{frame}

% %%%%%%%%%%%%%%%%%%%%%%%%%%%%%%%%%%%%%%%%%%%%%%%%%%%%%%%%%%%%%%%%%%%%%%%%%%%%%%%%%%
% \begin{frame}[fragile]\frametitle{Commands}
  % \begin{itemize}
    % \item \lstinline|/init| Create or update AGENTS.md file
	% \item \lstinline|/exit| Exit OpenCode. Aliases: \lstinline|/quit, /q|
	% \item \lstinline|/compact| Compact the current session. Alias: \lstinline|/summarize|
	% \item \lstinline|/help| Show the help dialog.
    % \item \lstinline|/undo or /redo|	Undo/Redo the last change
    % \item \lstinline|/models|	List available models
    % \item \lstinline|@|	Fuzzy search for files in your project
    % \item \lstinline|Tab|	Switch between Plan and Build mode
    % \item Create markdown files in the commands/ directory to define custom commands.
	% \item Create .opencode/commands/test.md and use it like \lstinline|/test|
  % \end{itemize}
 % \begin{lstlisting}
% ---
% description: Run tests with coverage
% agent: build
% model: anthropic/claude-3-5-sonnet-20241022
% ---

% Run the full test suite with coverage report and show any failures.
% Focus on the failing tests and suggest fixes. 
 % \end{lstlisting}
% \end{frame}

%%%%%%%%%%%%%%%%%%%%%%%%%%%%%%%%%%%%%%%%%%%%%%%%%%%%%%%%%%%%%%%%%%%%%%%%%%%%%%%%%%
\begin{frame}[fragile]\frametitle{MCP Servers:  External Tool Integration}
  \begin{itemize}
    \item \textbf{What:} Model Context Protocol (MCP) servers expose external tools and data sources to OpenCode. Any MCP-compatible server becomes a tool the agent can call natively.
    \item \textbf{Why:} Extends OpenCode beyond the filesystem:  databases, APIs, cloud services, monitoring tools, internal systems.
    \item \textbf{Popular MCP servers:} GitHub, GitLab, Slack, PostgreSQL, Puppeteer, Brave Search, Filesystem (extended), AWS, Docker, Linear, Jira, Notion, Sentry.
    \item \textbf{Configure in \texttt{opencode.json}:}
    \item \textbf{Tool naming:} MCP tools are prefixed by server name: \texttt{github\_create\_issue}, \texttt{postgres\_query}, etc. Use \texttt{mymcp\_*} in permissions to control all tools from one server.
    \item \textbf{Org defaults:} Deploy company-wide MCP servers via \texttt{.well-known/opencode}:  auto-loaded for all authenticated users.
    \item \textbf{Build your own:} Any HTTP/stdio process implementing the MCP spec becomes a server:  SDKs available in TypeScript, Python, Go.
  \end{itemize}
  
\end{frame}

%%%%%%%%%%%%%%%%%%%%%%%%%%%%%%%%%%%%%%%%%%%%%%%%%%%%%%%%%%%%%%%%%%%%%%%%%%%%%%%%%%
\begin{frame}[fragile]\frametitle{MCP Servers:  External Tool Integration}

\begin{lstlisting}
{ "mcp": {
    "github": {
      "type": "stdio",
      "command": "npx",
      "args": ["-y", "@modelcontextprotocol/server-github"],
      "env": { "GITHUB_TOKEN": "$GITHUB_TOKEN" }
    },
    "postgres": {
      "type": "stdio",
      "command": "npx",
      "args": ["-y", "@modelcontextprotocol/server-postgres",
               "postgresql://localhost/mydb"]
    }
  }
}
\end{lstlisting}  
\end{frame}

%%%%%%%%%%%%%%%%%%%%%%%%%%%%%%%%%%%%%%%%%%%%%%%%%%%%%%%%%%%%%%%%%%%%%%%%%%%%%%%%%%%
\begin{frame}[fragile]\frametitle{Custom Tools:  Extending the Agent's Hands}
  \begin{itemize}
    \item \textbf{What:} TypeScript/JS files that define new functions the LLM can call. Can wrap any language internally (Python, Bash, Rust, etc.).
    \item \textbf{Placement:} \texttt{.opencode/tools/<name>.ts} (project) or \texttt{\textasciitilde/.config/opencode/tools/<name>.ts} (global).
    \item \textbf{One file = one or more tools:} Single default export → tool named after file. Multiple named exports → tools named \texttt{<file>\_<export>}.
    \item \textbf{Multi-tool file example:}
    \item \textbf{Context available:} \texttt{context.agent}, \texttt{context.session}, \texttt{context.directory} (working dir), \texttt{context.worktree} (git root).
    \item \textbf{Override built-ins:} Same filename as a built-in (e.g., \texttt{bash.ts}) replaces the built-in completely:  useful for sandboxing or auditing.
    \item \textbf{Any language via shell:} Call Python/Bash scripts inside \texttt{execute()} using \texttt{child\_process}:  tool definition is TS, implementation is anything.
    \item \textbf{Schema validation:} Use \texttt{tool.schema} (Zod) for type-safe arg definitions with \texttt{.describe()} for LLM-facing documentation.
  \end{itemize}

\end{frame}



%%%%%%%%%%%%%%%%%%%%%%%%%%%%%%%%%%%%%%%%%%%%%%%%%%%%%%%%%%%%%%%%%%%%%%%%%%%%%%%%%%%
\begin{frame}[fragile]\frametitle{Custom Tools:  Extending the Agent's Hands}

\begin{lstlisting}
// .opencode/tools/math.ts
import { tool } from "@opencode-ai/plugin"
export const add = tool({ description: "Add two numbers",
  args: { a: tool.schema.number(), b: tool.schema.number() },
  async execute({ a, b }) { return a + b }
})
export const multiply = tool({ description: "Multiply",
  args: { a: tool.schema.number(), b: tool.schema.number() },
  async execute({ a, b }) { return a * b }
})
// creates: math_add, math_multiply
\end{lstlisting}  
\end{frame}


%%%%%%%%%%%%%%%%%%%%%%%%%%%%%%%%%%%%%%%%%%%%%%%%%%%%%%%%%%%%%%%%%%%%%%%%%%%%%%%%%
\begin{frame}[fragile]\frametitle{Models:  Provider \& Model Configuration}
  \begin{itemize}
    \item \textbf{75+ models supported} across: Anthropic, OpenAI, Google Gemini, Groq, Mistral, Cohere, xAI, DeepSeek, Qwen, Ollama (local), Hugging Face, AWS Bedrock, Azure, Vertex AI, and more.
	\item Chose models with long context window and have access to tools (function calling)
    \item \textbf{Switch models} in TUI: \texttt{/models} command → browse and select. Or \texttt{-m provider/model} CLI flag at launch.
    \item \textbf{Configure in \texttt{opencode.json}:}
    \item \textbf{Variants:} Define multiple configurations of the same model (different temperature, max\_tokens, reasoning effort) and cycle with \texttt{variant\_cycle} keybind.
    \item \textbf{Per-agent model:} Each custom agent can specify its own model:  e.g., \texttt{Plan} uses a cheap fast model; \texttt{Build} uses full Sonnet/Opus.
    \item \textbf{Small model for meta-tasks:} OpenCode uses a lightweight model (default: \texttt{gpt-5-nano} via Zen) for session titling and other housekeeping:  configurable separately.
    \item \textbf{Authenticate:} \texttt{/connect} in TUI → select provider → paste API key. Supports OAuth for GitHub/GitLab providers.
  \end{itemize}

\end{frame}

%%%%%%%%%%%%%%%%%%%%%%%%%%%%%%%%%%%%%%%%%%%%%%%%%%%%%%%%%%%%%%%%%%%%%%%%%%%%%%%%%
\begin{frame}[fragile]\frametitle{Models:  Provider \& Model Configuration}

\begin{lstlisting}
{ "model": "anthropic/claude-sonnet-4-6",
  "provider": {
    "anthropic": { "apiKey": "$ANTHROPIC_API_KEY" },
    "openai":    { "apiKey": "$OPENAI_API_KEY" },
    "ollama":    { "url": "http://localhost:11434" }
  }
}
\end{lstlisting}  
\end{frame}




%%%%%%%%%%%%%%%%%%%%%%%%%%%%%%%%%%%%%%%%%%%%%%%%%%%%%%%%%%%
%% USE CASES
%%%%%%%%%%%%%%%%%%%%%%%%%%%%%%%%%%%%%%%%%%%%%%%%%%%%%%%%%%%
\begin{frame}[fragile]\frametitle{Real-World Use Cases}
  \begin{itemize}
    \item \textbf{Understanding a Codebase:} \lstinline|You: Explain the authentication flow in this project| OpenCode will analyze your files and provide a detailed explanation of how authentication works.
  
    \item \textbf{Adding a Feature:} \lstinline|You: Add a dark mode toggle to the settings page| OpenCode will: Find the settings page, Identify the theme system, Implement the toggle, Update necessary styles

    \item \textbf{Debugging:} \lstinline|You: The login form is not submitting. Here's the error: [paste error]| OpenCode will trace the issue and suggest fixes.
	
    \item \textbf{Multi-File Refactoring:} Rename, restructure, or migrate patterns across the entire codebase atomically.
    \item \textbf{Documentation Generation:} Auto-generate README, API docs, and inline comments from code context.
    \item \textbf{Test Generation:} Scaffold unit/integration tests with awareness of existing test patterns.
    \item \textbf{Security Auditing:} Identify unsanitized inputs, SQL injection points, or insecure defaults.
  \end{itemize}
\end{frame}


%%%%%%%%%%%%%%%%%%%%%%%%%%%%%%%%%%%%%%%%%%%%%%%%%%%%%%%%%%%%%%%%%%%%%%%%%%%%%%%%%%
\begin{frame}[fragile]\frametitle{}

\begin{center}
{\Large Demo}
\end{center}
\end{frame}

%%%%%%%%%%%%%%%%%%%%%%%%%%%%%%%%%%%%%%%%%%%%%%%%%%%%%%%%%%%%%%%%%%%%%%%%%%%%%%%%%%
\begin{frame}[fragile]\frametitle{Initialize \& Understand}
  \begin{itemize}
    \item \textbf{Goal:} Show how OpenCode scans and understands a new project from scratch.
    \item \textbf{Setup:} Open terminal inside any Git repository.
    \item \textbf{Launch:} \lstinline|$ opencode|
    \item \textbf{Command:} \lstinline|> /init|
    \item \textbf{What Happens:} OpenCode scans all local files, infers stack (languages, deps, structure), and writes \texttt{AGENTS.md},  the AI's instruction manual for this project.
    \item \textbf{Key Point:} Every subsequent interaction is grounded in this context,  no hallucinated file paths.
  \end{itemize}
\end{frame}

%%%%%%%%%%%%%%%%%%%%%%%%%%%%%%%%%%%%%%%%%%%%%%%%%%%%%%%%%%%%%%%%%%%%%%%%%%%%%%%%%%
\begin{frame}[fragile]\frametitle{Fix \& Refactor}
  \begin{itemize}
    \item \textbf{Goal:} Demonstrate AI-driven code modification on a real file.
    \item \textbf{Command:} \lstinline|> @file:main.go fix the bug where user input isn't sanitized|
    \item \textbf{Agent Action:} Reads \texttt{main.go}, identifies the unsanitized input vector, proposes a targeted patch.
    \item \textbf{Review Flow:}
      \begin{itemize}
        \item Press \texttt{Tab} → switch to \textbf{Plan} mode: inspect the proposed diff before anything touches disk.
        \item Press \texttt{Tab} → switch to \textbf{Build} mode: apply the change.
      \end{itemize}
    \item \textbf{Key Point:} You stay in control,  the agent proposes, you decide.
  \end{itemize}
\end{frame}


%%%%%%%%%%%%%%%%%%%%%%%%%%%%%%%%%%%%%%%%%%%%%%%%%%%%%%%%%%%%%%%%%%%%%%%%%%%%%%%%%%
\begin{frame}[fragile]\frametitle{}

\begin{center}
{\Large Tips}
\end{center}
\end{frame}

%%%%%%%%%%%%%%%%%%%%%%%%%%%%%%%%%%%%%%%%%%%%%%%%%%%%%%%%%%%%%%%%%%%%%%%%%%%%%%%%%%
\begin{frame}[fragile]\frametitle{Tips for Beginners}
  \begin{itemize}
    \item Be Specific: 
Don't say "Make it better" but say, "Refactor the handleSubmit function to use async/await and add error handling"

    \item Provide Context: Don't say "Fix the bug" but say, "The login form throws 'undefined is not a function' when clicking submit. The error occurs in auth.ts line 45"

    \item Use Plan Mode First: For complex features, ask OpenCode to create a plan before implementation.

    \item Leverage Image Support: Drag and drop screenshots or mockups into the terminal—OpenCode can understand images!

    \item Create Project-Specific Configs: Different projects may need different settings. Use local opencode.json files.
  \end{itemize}
\end{frame}

%%%%%%%%%%%%%%%%%%%%%%%%%%%%%%%%%%%%%%%%%%%%%%%%%%%%%%%%%%%%%%%%%%%%%%%%%%%%%%%%%%
\begin{frame}[fragile]\frametitle{1. Root Rules File is Execution Memory}
  \begin{itemize}
    \item Your \texttt{OPENCODE.md} (or workspace rules file) should define:
    \begin{itemize}
        \item WHY — product purpose
        \item WHAT — repo map
        \item HOW — non-negotiable rules
    \end{itemize}
    \item Keep it short and authoritative.
    \item If it becomes a knowledge dump, OpenCode will ignore critical constraints.
    \item Clear, structural guidance beats long explanations.
  \end{itemize}
\end{frame}

%%%%%%%%%%%%%%%%%%%%%%%%%%%%%%%%%%%%%%%%%%%%%%%%%%%%%%%%%%%%%%%%%%%%%%%%%%%%%%%%%%
\begin{frame}[fragile]\frametitle{2. Turn Repeated Workflows into Templates}
  \begin{itemize}
    \item If you repeat instructions twice, template them.
    \item Create reusable folders such as:
  \end{itemize}

\begin{lstlisting}
.opencode/tasks/
.opencode/templates/
\end{lstlisting}

  \begin{itemize}
    \item Standardize code review, refactoring, and release workflows.
    \item Prevent prompt drift across sessions and teammates.
  \end{itemize}
\end{frame}

%%%%%%%%%%%%%%%%%%%%%%%%%%%%%%%%%%%%%%%%%%%%%%%%%%%%%%%%%%%%%%%%%%%%%%%%%%%%%%%%%%
\begin{frame}[fragile]\frametitle{3. Use Automation for Deterministic Guardrails}
  \begin{itemize}
    \item Memory is probabilistic. Tooling is deterministic.
    \item Use project-level commands, pre-commit hooks, and CI checks.
    \item Examples:
    \begin{itemize}
        \item Auto-run formatters
        \item Run tests after core changes
        \item Block unsafe directory edits (auth, migrations)
    \end{itemize}
    \item Deterministic rules belong in tooling, not prompts.
  \end{itemize}
\end{frame}

%%%%%%%%%%%%%%%%%%%%%%%%%%%%%%%%%%%%%%%%%%%%%%%%%%%%%%%%%%%%%%%%%%%%%%%%%%%%%%%%%%
\begin{frame}[fragile]\frametitle{4. Use docs/ for Progressive Disclosure}
  \begin{itemize}
    \item Keep architecture, ADRs, and runbooks in \texttt{docs/}.
    \item Reference documentation — do not paste it into context.
  \end{itemize}

\begin{lstlisting}
docs/
  architecture.md
  decisions/
  runbooks/
\end{lstlisting}

  \begin{itemize}
    \item Smaller context improves reliability.
  \end{itemize}
\end{frame}

%%%%%%%%%%%%%%%%%%%%%%%%%%%%%%%%%%%%%%%%%%%%%%%%%%%%%%%%%%%%%%%%%%%%%%%%%%%%%%%%%%
\begin{frame}[fragile]\frametitle{5. Add Local Rule Files Near Sharp Edges}
  \begin{itemize}
    \item Place rule files near high-risk directories.
    \item Document invariants, constraints, and warnings.
  \end{itemize}

\begin{lstlisting}
src/auth/RULES.md
src/persistence/RULES.md
infra/RULES.md
\end{lstlisting}

  \begin{itemize}
    \item Context is loaded exactly where it matters.
  \end{itemize}
\end{frame}

%%%%%%%%%%%%%%%%%%%%%%%%%%%%%%%%%%%%%%%%%%%%%%%%%%%%%%%%%%%%%%%%%%%%%%%%%%%%%%%%%%
\begin{frame}[fragile]\frametitle{6. Structure for Machine Onboarding}
  \begin{itemize}
    \item If a new engineer would be confused, OpenCode will be too.
    \item Use clear boundaries:
    \begin{itemize}
        \item \texttt{src/}
        \item \texttt{api/}
        \item \texttt{domain/}
        \item \texttt{infra/}
    \end{itemize}
    \item Keep naming consistent and nesting shallow.
  \end{itemize}
\end{frame}

%%%%%%%%%%%%%%%%%%%%%%%%%%%%%%%%%%%%%%%%%%%%%%%%%%%%%%%%%%%%%%%%%%%%%%%%%%%%%%%%%%
\begin{frame}[fragile]\frametitle{7. Smaller Context Beats Larger Context}
  \begin{itemize}
    \item More files in context does not mean better output.
    \item Large context dilutes signal and increases hallucination risk.
    \item Provide:
    \begin{itemize}
        \item Exact files
        \item Exact constraints
        \item Exact task definition
    \end{itemize}
    \item Precision beats volume.
  \end{itemize}
\end{frame}

%%%%%%%%%%%%%%%%%%%%%%%%%%%%%%%%%%%%%%%%%%%%%%%%%%%%%%%%%%%%%%%%%%%%%%%%%%%%%%%%%%
\begin{frame}[fragile]\frametitle{8. Document Architecture Decisions (ADRs)}
  \begin{itemize}
    \item OpenCode cannot infer historical decisions.
    \item Document why technologies and boundaries were chosen.
  \end{itemize}

\begin{lstlisting}
docs/decisions/
\end{lstlisting}

  \begin{itemize}
    \item Prevent the AI from “optimizing” your architecture away.
  \end{itemize}
\end{frame}

%%%%%%%%%%%%%%%%%%%%%%%%%%%%%%%%%%%%%%%%%%%%%%%%%%%%%%%%%%%%%%%%%%%%%%%%%%%%%%%%%%
\begin{frame}[fragile]\frametitle{9. Treat Prompts as Versioned Assets}
  \begin{itemize}
    \item Store reusable prompts in your repository.
  \end{itemize}

\begin{lstlisting}
tools/prompts/
feature-implementation.md
refactor-checklist.md
bug-triage.md
\end{lstlisting}

  \begin{itemize}
    \item Version and improve them over time.
    \item Avoid copy-paste drift across sessions.
  \end{itemize}
\end{frame}

%%%%%%%%%%%%%%%%%%%%%%%%%%%%%%%%%%%%%%%%%%%%%%%%%%%%%%%%%%%%%%%%%%%%%%%%%%%%%%%%%%
\begin{frame}[fragile]\frametitle{10. Structure > Prompting Skill}
  \begin{itemize}
    \item Strong structure transforms OpenCode into a workspace-native engineer.
    \item Combine:
    \begin{itemize}
        \item File-based rules
        \item Workflow templates
        \item Automated guardrails
        \item Documented architecture
        \item Minimal scoped context
    \end{itemize}
    \item Prompts are temporary. Structure is permanent.
  \end{itemize}
\end{frame}

%%%%%%%%%%%%%%%%%%%%%%%%%%%%%%%%%%%%%%%%%%%%%%%%%%%%%%%%%%%%%%%%%%%%%%%%%%%%%%%%%%
\begin{frame}[fragile]\frametitle{}

\begin{center}
{\Large Summary}
\end{center}
\end{frame}

%%%%%%%%%%%%%%%%%%%%%%%%%%%%%%%%%%%%%%%%%%%%%%%%%%%%%%%%%%%
%% WHY OPENCODE
%%%%%%%%%%%%%%%%%%%%%%%%%%%%%%%%%%%%%%%%%%%%%%%%%%%%%%%%%%%
\begin{frame}[fragile]\frametitle{Why Use OpenCode?}
  \begin{itemize}
    \item \textbf{Privacy-First:} Code never leaves your machine by default; no telemetry, no vendor data retention.
    \item \textbf{Cost Efficiency:} No subscription fees,  pay only for actual API token usage; local models are free.
    \item \textbf{No Vendor Lock-in:} Switch AI providers in seconds via config; one tool, all models.
    \item \textbf{Terminal-Native Speed:} Faster context-switch than IDE plugins; integrates directly into your existing shell workflow.
    \item \textbf{Transparent \& Auditable:} Fully open-source (Apache 2.0),  inspect, fork, or extend anything.
    \item \textbf{Enterprise-Ready:} Supports self-hosted LLMs and air-gapped (physically isolated from all other networks) environments.
    \item \textbf{Bottom Line:} Faster dev loop, zero lock-in, full code privacy,  AI that works \textit{with} your terminal.	
  \end{itemize}
\end{frame}

%%%%%%%%%%%%%%%%%%%%%%%%%%%%%%%%%%%%%%%%%%%%%%%%%%%%%%%%%%%%%%%%%%%%%%%%%%%%%%
\begin{frame}[fragile]\frametitle{Quick Reference:  Config Cheat Sheet}
  \begin{itemize}
    \item \textbf{File locations} (\texttt{P} = project \texttt{.opencode/}, \texttt{G} = global \texttt{\textasciitilde/.config/opencode/}):
      \begin{itemize}
        \item Rules: \texttt{P: AGENTS.md} · \texttt{G: AGENTS.md}
        \item Agents: \texttt{P: agents/<name>.md} · \texttt{G: agents/<name>.md}
        \item Skills: \texttt{P: skills/<name>/SKILL.md} · \texttt{G: skills/<name>/SKILL.md}
        \item Tools: \texttt{P: tools/<name>.ts} · \texttt{G: tools/<name>.ts}
        \item Modes: \texttt{P: modes/<name>.md} · \texttt{G: modes/<name>.md}
        \item Commands: \texttt{opencode.json → "command": \{...\}}
        \item MCP: \texttt{opencode.json → "mcp": \{...\}}
        \item Permissions: \texttt{opencode.json → "permission": \{...\}}
      \end{itemize}
    \item \textbf{Key CLI commands:}
      \begin{itemize}
        \item \texttt{opencode agents create}:  scaffold a new agent interactively
        \item \texttt{/init}:  generate \texttt{AGENTS.md} for current project
        \item \texttt{/models}:  browse and switch models
        \item \texttt{/connect}:  authenticate a new provider
        \item \texttt{/share}:  share current session URL
      \end{itemize}
    \item \textbf{Priority order} (highest → lowest): CLI flags → \texttt{OPENCODE\_CONFIG\_DIR} → \texttt{.opencode/opencode.json} → \texttt{\textasciitilde/.config/opencode/opencode.json} → \texttt{.well-known/opencode} (org defaults).
    \item \textbf{Docs:} \href{https://opencode.ai/docs}{\texttt{opencode.ai/docs}} · \textbf{GitHub:} \href{https://github.com/anomalyco/opencode}{\texttt{github.com/anomalyco/opencode}}
  \end{itemize}
\end{frame}

%%%%%%%%%%%%%%%%%%%%%%%%%%%%%%%%%%%%%%%%%%%%%%%%%%%%%%%%%%%
%% References
%%%%%%%%%%%%%%%%%%%%%%%%%%%%%%%%%%%%%%%%%%%%%%%%%%%%%%%%%%%
\begin{frame}[fragile]\frametitle{ Resources}
  \begin{itemize}
	\item GitHub: \href{https://github.com/opencode-ai/opencode}{\texttt{github.com/opencode-ai/opencode}}
	\item Docs: \href{https://opencode.ai/docs}{\texttt{opencode.ai/docs}}
	\item NXCode.io ``OpenCode Tutorial 2026: Install \& Setup After Anthropic Block (Updated Guide)''
	\item JP Caparas, ``The definitive guide to OpenCode: from first install to production workflows'', Dev Genius / reading.sh, Feb 2026. \href{https://blog.devgenius.io/the-definitive-guide-to-opencode-from-first-install-to-production-workflows-aae1e95855fb}{\texttt{blog.devgenius.io}}
  \end{itemize}
\end{frame}
